\section{Self-Contained Proofs For The Exhaustion Chain}
\label{sec:appendix-c-proofs}

This appendix gives the full proofs of the capacity-exhaustion steps
summarized in Section~6, so the article is self-contained. Each entry states the
property's population, scope, derivation, and boundary. Its proofs are the
authority for the exhaustion claims made in this article; the wider research
records use the same notation and provide supplementary context.

Notation shared across the entries: $D=Q^2-Q-3$, $a_i=1-2/r_i$,
$P_i=\prod_{j<i}a_j$, $P_m=\prod_{j<m}a_j$, $T=N_0P_m$, and
$W_-=\sum_{i<m}P_m/P_i$. The per-layer squared harmful-excess envelope is
written $X_i$ throughout (in earlier research records the capacity
theorem writes it $M_i$; the native-period hybrid theorem renamed it $X_i$
to avoid collision with the native modulus $M_k$, and this article follows
the later convention).

\subsection{Sharp Harmful-Capacity Excess Envelope}
\label{sec:appendix-c-1}

Fix one filter layer's residue histogram, restricted to the two harmful
classes: for every incoming prime $r>2$, every common residue-class capacity
$B\ge0$, and every population $0\le N\le rB$, let the residue counts $c_a$
($a\bmod r$) satisfy $0\le c_a\le B$ and $\sum c_a=N$. The two harmful
classes hold $N_{\text{harm}}=c_0+c_{-2}$ starts, and the signed harmful
excess is $b=N_{\text{harm}}-2N/r$.

The total $N_{\text{harm}}$ is constrained by two facts. First, each harmful
class holds at most $B$, so $N_{\text{harm}}\le 2B$. Also
$N_{\text{harm}}\le N$. Second, the other $r-2$ classes hold at most
$(r-2)B$ of the $N$ starts, forcing at least $N-(r-2)B$ into the harmful
pair. The sixfold-capacity envelope proves both endpoints are attainable,
giving the exact feasible interval
\begin{equation*}
\ell\le N_{\text{harm}}\le u,
\qquad
\ell=\max(0,N-(r-2)B),
\qquad
u=\min(N,2B).
\end{equation*}

Because $b^2=(N_{\text{harm}}-2N/r)^2$ is convex in $N_{\text{harm}}$, its
maximum over $[\ell,u]$ occurs
at an endpoint:
\begin{equation*}
b^2
\le
X_{r,N,B}
:=
\max\left\{
\left(\ell-\frac{2N}{r}\right)^2,
\left(u-\frac{2N}{r}\right)^2
\right\}.
\end{equation*}

Both endpoints are attainable, so this bound is sharp. The conditioned-chain
upper envelope follows by summing the per-layer sharp bounds with the energy
coefficients $\alpha_i=w_ir_i/[2(r_i-2)]$:
\begin{equation*}
E_b\le\mathcal U_{\mathrm{cap}}:=\sum_i\alpha_iX_i
\qquad\blacksquare\ \text{[Q.E.D.]}
\end{equation*}

The bound is sharp one layer at a time but need not be sharp
over a chain, because the histograms attaining each $X_i$ separately may not
co-arise from one nested survivor sequence. A cross-layer CRT restriction
could lower the true aggregate energy below $\mathcal U_{\mathrm{cap}}$. The
property does not prove $\mathcal U_{\mathrm{cap}}\lt T^2/(2W_-)+\Gamma_{\mathrm{cap}}$;
it reduces the capacity-only route to that explicit inequality.

\subsection{Capacity-Envelope Width Floor Needs Population Slack}
\label{sec:appendix-c-2}

For every $r\ge5$, $B\ge0$, $0\le N\le rB$, over one filter layer's feasible
harmful-count interval, this property extracts the explicit lower bound on
$X_{r,N,B}$ supplied by
the \emph{width} of the feasible interval $[\ell,u]$. Write
$c=(\ell+u)/2$, $h=(u-\ell)/2$. The farther endpoint from $2N/r$ has distance
\begin{equation*}
\max\left(\left|\frac{2N}{r}-(c-h)\right|,\left|\frac{2N}{r}-(c+h)\right|\right)
=h+\left|\frac{2N}{r}-c\right|\ge h.
\end{equation*}

Therefore
\begin{equation*}
X_{r,N,B}\ge\frac{(u-\ell)^2}{4}.
\end{equation*}

The width has the closed form
\begin{equation*}
u-\ell=\min(N,2B,rB-N).
\end{equation*}

This follows by splitting the feasible range into three parts. If
$0\le N\le 2B$, then $u=N$, $\ell=0$, and $u-\ell=N$. If
$2B\le N\le(r-2)B$, then $u=2B$, $\ell=0$, and $u-\ell=2B$. If
$(r-2)B\le N\le rB$, then $u=2B$, $\ell=N-(r-2)B$, and $u-\ell=rB-N$. The
formulas agree at the shared endpoints, proving the minimum formula.
Combining,
\begin{equation*}
X_{r,N,B}\ge\frac14\min(N,2B,rB-N)^2
\qquad\blacksquare\ \text{[Q.E.D.]}
\end{equation*}

\textbf{Zero characterization.} Because $X_{r,N,B}$ is the maximum of two squares,
$X_{r,N,B}=0$ iff both endpoints equal $2N/r$, which requires $u=\ell$. By the
width formula this is equivalent to $\min(N,2B,rB-N)=0$. Since $2B>0$ and
$0\le N\le rB$, this occurs exactly at
\begin{equation*}
X_{r,N,B}=0\iff N\in\{0,rB\}.
\end{equation*}

This is the obstruction: the envelope vanishes at both the empty and the
fully-occupied populations. Theorem using only $r$ and $B$ cannot force a
positive floor, because the fully-occupied profile $N=rB$ is positive and has
zero capacity envelope.

\subsection{Fixed Seven Cut Cannot Clear The Original Threshold}
\label{sec:appendix-c-3}

Consider the cut immediately after filter $7$, so $k=2$, for every chain
$r_0=5,r_1=7,r_2=11,\ldots$ with $Q\ge17$ and $m\ge37$, examining the suffix
capacity term at the first untouched layer (filter $11$) of a conditioned
chain and assuming the seven-layer density floor's local-count threshold at
filter $11$. The native-period
hybrid envelope leaves every coordinate $i\ge2$ under its separate capacity
bound, giving
\begin{equation*}
\mathcal U_2^{\mathrm{hyb}}\ge\alpha_2X_2.
\end{equation*}

The first three multiplicative factors are $a_0=3/5$, $a_1=5/7$, $a_2=9/11$,
so $P_3=a_0a_1a_2=27/77$. Since $w_2=P_m/P_3$ and $\alpha_2=w_2/(2a_2)$,
\begin{equation*}
\alpha_2=\frac{P_m}{2a_2P_3}=\frac{P_m}{2\cdot(9/11)\cdot(27/77)}
=\frac{847}{486}P_m.
\end{equation*}

The filter-$11$ capacity floor gives $X_2\ge B_{11}^2$ with
$B_{11}=\lfloor D/66\rfloor+1\ge D/66$. Therefore
\begin{equation*}
\mathcal U_2^{\mathrm{hyb}}\ge\frac{847}{486}P_m\left(\frac{D}{66}\right)^2.
\end{equation*}

For the threshold upper bound, $T=N_0P_m$ and $W_-=\sum_{i<m}P_m/P_i\ge mP_m$
(since $P_i\le1$). Before filter $5$ every 2-gap start is $5\bmod6$, so
$N_0\le\lfloor D/6\rfloor+1\le D/6+1$. Hence
\begin{equation*}
\frac{T^2}{2W_-}\le\frac{P_m}{2m}\left(\frac D6+1\right)^2.
\end{equation*}

The suffix exceeds the threshold whenever
\begin{equation*}
\frac{847}{486}\left(\frac{D}{66}\right)^2>\frac1{2m}\left(\frac D6+1\right)^2,
\end{equation*}

equivalently $m>\frac{29403}{847}(1+6/D)^2$. For $Q\ge17$, $D\ge269$, so the
right side is at most $\frac{29403}{847}(275/269)^2$. The integer comparison
\begin{equation*}
29403\cdot275^2=2{,}223{,}601{,}875\lt2{,}267{,}721{,}379=37\cdot847\cdot269^2
\end{equation*}

shows this is strictly below $37$. Therefore $m\ge37$ proves
\begin{equation*}
\mathcal U_2^{\mathrm{hyb}}>\frac{T^2}{2W_-}
\qquad\blacksquare\ \text{[Q.E.D.]}
\end{equation*}

This proves the \emph{envelope} cannot certify survival through this
fixed cut. It does not bound the actual energy $E_b$. It does not address
later optimized cuts ($k\ge3$), the capacity-relaxed threshold, or localized
residue information.

\subsection{Every Fixed Native Cut Fails The Original Threshold}
\label{sec:appendix-c-4}

Generalize C.3 to an arbitrary cut $k$, for every chain
$5\le r_0<\cdots<r_{m-1}<Q$ with $Q\ge17$ and $2\le k<m$, over the suffix
capacity term at the first untouched layer of an arbitrary native cut. The
seven-layer density floor at layer $r_k$ gives
$X_k\ge B_k^2$ with $B_k=\lfloor D/(6r_k)\rfloor+1\ge D/(6r_k)$, and
$\alpha_k=P_m/(2P_ka_k^2)$. Therefore
\begin{equation*}
\mathcal U_k^{\mathrm{hyb}}
\ge\frac{P_mD^2}{72P_ka_k^2r_k^2}.
\end{equation*}

The threshold upper bound is unchanged:
$T^2/(2W_-)\le(P_m/2m)(D/6+1)^2$. The suffix exceeds the threshold whenever
\begin{equation*}
\frac{P_mD^2}{72P_ka_k^2r_k^2}>\frac{P_m}{2m}\left(\frac D6+1\right)^2.
\end{equation*}

Canceling $P_m>0$ and rearranging gives $m>P_ka_k^2r_k^2(1+6/D)^2$. Since
$a_kr_k=r_k-2$,
\begin{equation*}
m>P_k(r_k-2)^2\left(1+\frac6D\right)^2
\quad\Longrightarrow\quad
\mathcal U_k^{\mathrm{hyb}}>\frac{T^2}{2W_-}.
\end{equation*}

Every fixed cut eventually fails: for fixed $k$, $P_k$ and $r_k$ are constants
while $(1+6/D)^2\to1$ as $Q$ grows. Along any family with $m\to\infty$, the
fixed cut violates the necessary condition. Thus a cut capable of clearing the
original threshold must move: $k=k(Q)\to\infty$.

A parameter-free lower bound on the cut prime follows from $k\ge2$, giving
$P_k\le P_2=(3/5)(5/7)=3/7$. Threshold clearance would require
\begin{equation*}
m\le\frac37(r_k-2)^2\left(1+\frac6D\right)^2,
\end{equation*}

equivalently
\begin{equation*}
\mathcal U_k^{\mathrm{hyb}}\lt\frac{T^2}{2W_-}
\quad\Longrightarrow\quad
r_k\ge2+\frac{\sqrt{7m/3}}{1+6/D}
\qquad\blacksquare\ \text{[Q.E.D.]}
\end{equation*}

\textbf{Recovery of C.3.} For the cut after filter $7$, $k=2$, $r_2=11$,
$P_2=3/7$, and $P_2(r_2-2)^2=(3/7)\cdot81=243/7=29403/847$, recovering
C.3's constant exactly.

The bound $r_k\ge2+\sqrt{7m/3}/(1+6/D)$ uses no estimate for
the distribution of primes. It is a necessary condition on the cut prime, not
sufficient: it does not prove such an $r_k$ exists in the chain. Converting
the prime bound into a cut-index bound requires the prime number theorem
(Appendix~C.5).

\subsection{Moving Cut Loses Complete Native Blocks}
\label{sec:appendix-c-5}

For every chain $5\le r_i<Q$, cut $2\le k<m$, over a native cut that moves
outward with the future head together with the prime-counting function, the
exact logarithmic-squared inequality below holds under five stated finite
hypotheses. The asymptotic corollary uses Bertrand's postulate and the prime
number theorem as explicit external dependencies; the exact theorem itself is
finite, and only the asymptotic corollary is external.

The native modulus at cut $k$ is $M_k=\prod_{p<r_k}p=2\cdot3\prod_{i<k}r_i$.
Since $r_{k-1}$ is the prime immediately before $r_k$,
\begin{equation*}
\log M_k=\vartheta(r_{k-1}).
\end{equation*}

Assume the following five hypotheses:

\begin{enumerate}
\item the seven-layer density floor holds at layer $r_k$;
\item the cut clears the original threshold, $\mathcal U_k^{\mathrm{hyb}}\lt T^2/(2W_-)$;
\item the native modulus fits the interval, $M_k\le H$ where $H=D+1=Q^2-Q-2$;
\item for some constant $c>0$, $\vartheta(r_{k-1})\ge cr_{k-1}$; and
\item Bertrand's inequality $r_k\lt2r_{k-1}$.
\end{enumerate}

From C.4, hypothesis 2 forces
$r_k\ge2+\sqrt{7m/3}/(1+6/D)$. From hypotheses 3--5,
\begin{equation*}
\log H\ge\log M_k=\vartheta(r_{k-1})\ge cr_{k-1}>\frac c2 r_k.
\end{equation*}

Combining the lower and upper requirements on $r_k$ and rearranging,
\begin{equation*}
m\lt\frac37\left(1+\frac6D\right)^2\left(\frac{2\log H}{c}-2\right)^2.
\end{equation*}

This is the exact finite theorem: a threshold-clearing cut with at least one
complete native block forces the chain length to be at most $O(\log^2H)$.

\textbf{Prime-number-theorem corollary (external).} The asymptotics
$\vartheta(x)\sim x$ and $\pi(x)\sim x/\log x$ are external to the project's
verification. For the actual full chain, $m=\pi(Q)-3\sim Q/\log Q$, whereas
$\log^2H\sim4\log^2Q$. Therefore
\begin{equation*}
\frac{m}{\log^2H}\sim\frac{Q}{4\log^3Q}\longrightarrow\infty.
\end{equation*}

The exact logarithmic-squared necessary condition fails for all sufficiently
large $Q$. Hence, under the seven-layer density floor at the first suffix
layer,
\begin{equation*}
\mathcal U_k^{\mathrm{hyb}}\lt\frac{T^2}{2W_-}
\quad\Longrightarrow\quad
M_k>H
\end{equation*}

for every sufficiently large $Q$ and every cut $k$. There are no complete
native blocks to cancel.

When $M_k>H$, the native-period hybrid envelope still
constrains the single incomplete interval block. C.6 addresses whether that
constraint can supply the missing gain. The asymptotic conclusion depends
explicitly on PNT and Bertrand; the exact inequality remains valid without
them under the five stated hypotheses.

\subsection{Incomplete-Block Bessel Excludes No Capacity}
\label{sec:appendix-c-6}

Under the hypotheses of C.4--C.5, for every cut $2\le k<m$ with $M_k>H$, the
native-period hybrid envelope's interval remainder on the single incomplete
native block left is $s_k=H$: the whole window is one incomplete block (the
asymptotic scale uses Bertrand and PNT as external dependencies). This entry
proves the
normalized capacity box then fits inside that budget, so the overflow $e_k$
vanishes and the hybrid envelope collapses back to the all-capacity
envelope.

The native-period hybrid envelope's exact norm at coordinate $i<k$ is
\begin{equation*}
q_{i,k}=\frac{M_kP_i(r_i-2)}{3r_i^2}.
\end{equation*}

The capacity numerator obeys $X_i\le N_i^2$. Since conditioned populations
decrease, $N_i\le N_0$, and before filter $5$ every 2-gap start is $5\bmod6$,
so $N_0\le\lfloor D/6\rfloor+1\le D/5$ (using $D\ge269>30$). Therefore
\begin{equation*}
X_i\le\frac{D^2}{25}.
\end{equation*}

For the denominator, the function $x\mapsto(x-2)/x^2$ is decreasing for
$x\ge4$, so for $i<k$,
\begin{equation*}
\frac{r_i-2}{r_i^2}\ge\frac{r_k-2}{r_k^2},
\qquad
q_{i,k}\ge\frac{M_kP_k(r_k-2)}{3r_k^2}.
\end{equation*}

Summing the $k$ prefix coordinates,
\begin{equation*}
\sum_{i<k}\frac{X_i}{q_{i,k}}
\le\frac{3kD^2r_k^2}{25M_kP_k(r_k-2)}.
\end{equation*}

The overflow quantification defines
$e_k=(\sum_{i<k}X_i/q_{i,k}-s_k)_+$. When $M_k>H$, $s_k=H$. The normalized
sum is at most $H$---giving $e_k=0$---whenever
\begin{equation*}
M_kP_k\ge\frac{3kD^2r_k^2}{25H(r_k-2)}.
\end{equation*}

\textbf{Prime-number-theorem scale (external).} The prefix product satisfies
$M_kP_k=6\prod_{i<k}(r_i-2)\ge M_k/2^k$. Using PNT and Bertrand externally,
C.5 forces $r_k\gg\sqrt{Q/\log Q}$ for any potentially successful cut, hence
$\log(M_kP_k)\sim r_{k-1}\gg\sqrt{Q/\log Q}$. The logarithm of the right side
of the zero-overflow criterion is only $O(\log Q)$, since $k<Q$,
$D\lt H\lt Q^2$, and $r_k<Q$. The criterion therefore holds for every sufficiently
large $Q$, giving
\begin{equation*}
e_k=0,\qquad\mathcal U_k^{\mathrm{hyb}}=\mathcal U_{\mathrm{cap}}.
\end{equation*}

\textbf{Exhaustion of the original native hybrid.} Combining C.3--C.6: fixed cuts
are excluded by C.3--C.4; any potentially successful moving cut has
$M_k>H$ (C.5) and then $\mathcal U_k^{\mathrm{hyb}}=\mathcal U_{\mathrm{cap}}$
(C.6). Since C.3 gives $\mathcal U_{\mathrm{cap}}\ge\mathcal U_2^{\mathrm{hyb}}>T^2/(2W_-)$,
the capacity-plus-native-Bessel envelope satisfies
\begin{equation*}
\mathcal U_{\mathrm{hyb}}\ge\frac{T^2}{2W_-}
\end{equation*}

for every sufficiently large complete chain satisfying the seven-layer
density floor. The current capacity-plus-native-Bessel envelope cannot
certify the terminal survival threshold on an unbounded family.

This is a method obstruction, not a refutation of the
seven-layer density floor or the terminal survival estimate. It does not
address the capacity-relaxed threshold
$T^2/(2W_-)+\Gamma_{\mathrm{cap}}$ (handled by the stability-gap theorem,
summarized in Section~6.6) or a localized upper bound for the actual $E_b$.
